\documentclass{beamer}
\usepackage{showframe}

\usepackage[american]{babel}
\usepackage[T1]{fontenc}

\usepackage{fontspec}
\setsansfont{Jost}

\usepackage{emoji}
\setemojifont{Noto Color Emoji}

\usepackage{microtype}
\usepackage{csquotes}

\usepackage{xcolor}
\usepackage{hyperref}

\usepackage{amsmath}
\usepackage{amssymb}
\usepackage{amsthm}
\usepackage{mathtools}
\usepackage{stmaryrd}
\usepackage{wasysym}

\usepackage{tikz}
\usetikzlibrary{calc, fit, patterns, positioning, shapes, arrows.meta, overlay-beamer-styles}

\usepackage{ulem}
\usepackage{booktabs}
\usepackage{xparse}

% ==============================================================================
% Macros for Relational Separation Logic and RTL Compiler Verification
% ==============================================================================

% Operational Semantics and Reductions
\newcommand{\step}[3]{#1 \vdash #2 \twoheadrightarrow #3}
\newcommand{\steps}[3]{#1 \vdash #2 \twoheadrightarrow^{\ast} #3}
\newcommand{\psteps}[3]{#1 \vdash #2 \twoheadrightarrow^{+} #3}
\newcommand{\nstep}[4]{#1 \vdash #2 \underset{#3}{\twoheadrightarrow} #4}

\newcommand{\isval}[1]{\textsf{is\_value}\left(#1\right)}
\newcommand{\bfinal}[3]{\textsf{both\_final}_{#1}\left(#2, #3\right)}
\newcommand{\canp}[2]{\textsf{can\_progress}_{#1}\left(#2\right)}
\newcommand{\stuck}[2]{\textsf{stuck}_{#1}\left(#2\right)}
\newcommand{\Some}[1]{\textsf{Some}\left(#1\right)}
\newcommand{\None}{\textsf{None}}

% Behavioral Predicates
\newcommand{\Spec}{\textsf{Spec}}
\newcommand{\Safe}[1]{\textsf{Safe}(#1)}
\newcommand{\Term}[2]{\textsf{Terminating}(#1, #2)}
\newcommand{\Div}{\textsf{Diverging}}
\newcommand{\Undef}{\textsf{Undefined}}

% Data Structures
\newcommand{\List}[1]{\textsf{List}\left(#1\right)}
\newcommand{\Stream}[1]{\textsf{Stream}\left(#1\right)}

% Simulation Relations
\newcommand{\isim}[3]{#1 \lesssim #2 \; \{#3\}}
\newcommand{\fsim}[5]{#1 \leftindex[I]_{#2}\lesssim_{#3} #4 \; \{#5\}}
\newcommand{\incirc}[1]{\raisebox{.5pt}{\textcircled{\raisebox{-.9pt}{#1}}}}

% Value Types and Constructors
\newcommand{\val}{\textsf{val}}
\newcommand{\loc}{\textsf{loc}}
\newcommand{\VInt}[1]{\textsf{VInt}\left(#1\right)}
\newcommand{\VBool}[1]{\textsf{VBool}\left(#1\right)}
\newcommand{\VPtr}[1]{\textsf{VPtr}\left(#1\right)}
\newcommand{\VUndef}{\textsf{VUndef}}

% Memory Cell Types and Constructors
\newcommand{\cell}{\textsf{cell}}
\newcommand{\Allocated}[1]{\textsf{Allocated}\left(#1\right)}
\newcommand{\Freed}{\textsf{Freed}}
\newcommand{\memory}{\textsf{memory}}

% Relational Value Similarity
\newcommand{\sameval}[3]{#2 \mathrel{\bowtie}_{#1} #3}

% Regmap Operations
\newcommand{\getreg}[2]{\ensuremath{#2[#1]}}
\newcommand{\setreg}[3]{\ensuremath{#3[#1 \mapsto #2]}}

% Memory Operations
\newcommand{\getat}[2]{\ensuremath{#2[#1]}}
\newcommand{\setat}[3]{\ensuremath{#3[#1 \mapsto \Allocated{#2}]}}
\newcommand{\allocat}[3]{\ensuremath{#3 \uplus \{#1 \mapsto \Allocated{#2}\}}}
\newcommand{\freeat}[2]{\ensuremath{#2[#1 \mapsto \Freed]}}

% State and Stack Constructors
\newcommand{\State}[3]{\textsf{State}\left(#1, #2, #3\right)}
\newcommand{\CallState}[2]{\textsf{CallState}\left(#1, #2\right)}
\newcommand{\ReturnState}[1]{\textsf{ReturnState}\left(#1\right)}
\newcommand{\Stackframe}[4]{\textsf{Stackframe}\left(#1, #2, #3, #4\right)}
\newcommand{\pcstate}{\textsf{pcstate}}
\newcommand{\rtlstate}{\textsf{rtl\_state}}

% RTL Instruction Set (single-line representations)
\newcommand{\Inop}[1]{\mathtt{nop} \to #1}
\newcommand{\Iop}[4]{#1 \coloneqq #2\; #3 \to #4}
\newcommand{\Iload}[3]{#1 \coloneqq {!#2} \to #3}
\newcommand{\Istore}[3]{{!#1} \coloneqq #2 \to #3}
\newcommand{\Ialloc}[2]{#1 \coloneqq \mathtt{alloc}() \to #2}
\newcommand{\Ifree}[2]{\mathtt{free}\; #1 \to #2}
\newcommand{\Icond}[3]{\mathtt{if}\; #1\; \mathtt{then}\; #2\; \mathtt{else}\; #3}
\newcommand{\Icall}[4]{#1 \coloneqq \mathtt{call}\; #2(#3) \to #4}
\newcommand{\Iret}[1]{\mathtt{ret}\; #1}
\newcommand{\Ireturn}[1]{\mathtt{ret}\; #1}

% RTL Helpers
\newcommand{\evalop}[2]{\textsf{eval\_op}\left(#1, #2\right)}
\newcommand{\findfun}[2]{\textsf{find\_fun}\left(#1, #2\right)}
\newcommand{\initregs}[2]{\textsf{init\_regs}\left(#1, #2\right)}

% Custom Operators
\makeatletter
\newcommand{\slantsqsubseteq}{\mathrel{\mathpalette\draw@slantsqsubseteq\relax}}
\newcommand{\draw@slantsqsubseteq}[2]{%
  \vcenter{\hbox{%
    \begin{tikzpicture}[
      baseline=0,
      line width=0.04em,
      line cap=round,
      line join=round,
      scale=0.08em
    ]
      % top arm down to base
      \draw (0.7em, 0.8em) -- (0em, 0.33em) -- (0.7em, 0.33em);
      % Underbar
      \draw (0em, 0.1em) -- (0.7em, 0.1em);
    \end{tikzpicture}%
  }}%
}
\makeatother


\newcommand{\Prop}{\textsf{Prop}} % Maybe \mathbb{P} ?
\newcommand{\reg}{\textsf{reg}}
\newcommand{\mem}{\textsf{memory}}
\newcommand{\regbank}{\textsf{regbank}}
\newcommand{\listT}[1]{\textsf{list}~{#1}}
\newcommand{\node}{\textsf{node}}
\newcommand{\instr}{\textsf{instr}}
\newcommand{\ident}{\textsf{ident}}
\newcommand{\pcatis}[3]{#1[#2] = \langle#3\rangle}
\newcommand{\true}{\textsf{true}}
\newcommand{\false}{\textsf{false}}

% Relational Separation Logic Macros
\newcommand{\rProp}{\textsf{rProp}}
\newcommand{\purel}[1]{\ulcorner #1 \urcorner}
\newcommand{\emp}{\textsf{emp}}
\providecommand{\sep}{\mathbin{\ast}}
\renewcommand{\sep}{\mathbin{\ast}}
\newcommand{\wand}{\mathrel{-\mkern-6mu\ast}}
\newcommand{\ent}{\vdash}
\newcommand{\boxm}{\Box}
\newcommand{\ptrt}[2]{#1 \to_{\text{t}} #2}
\newcommand{\ptrs}[2]{#1 \to_{\text{s}} #2}
\newcommand{\freet}[1]{\textsf{free}_{\text{t}}\left(#1\right)}
\newcommand{\frees}[1]{\textsf{free}_{\text{s}}\left(#1\right)}
\newcommand{\meminj}[2]{\textsf{meminj}_{#2}\left(#1\right)}
\newcommand{\sameargs}[3]{\textsf{same\_args}\left(#1, #2, #3\right)}
\newcommand{\samer}[3]{\textsf{samer}\left(#1, #2, #3\right)}
\newcommand{\hoareq}[4]{\{#1\}\; #2 \lesssim #3\; \{#4\}}
\newcommand{\bigast}{\mathop{\raisebox{-2pt}{\Large$\ast$}}\displaylimits}


\usetheme{Singapore}

\colorlet{srccol}{orange!80!black}
\colorlet{tgtcol}{blue!80!black}

\makeatletter % https://tex.stackexchange.com/a/656319/16595
\tikzset{parse let/.code={\def\tikz@cc@stop@let in{}\tikz@let@command et #1in}}
\makeatother

\tikzset{
  bbox/.style={
    % draw, red
    draw=none
  },
  text width between/.style args={#1 and #2}{
    parse let={\p@=($(#2)-(#1)$)},
    text width/.expanded={abs(\x@)-2*(\noexpand\pgfkeysvalueof{/pgf/inner xsep})}
  }
}

\def\subnode#1#2%
{\tikz[remember picture,baseline=(#1.base),inner sep=0pt,
  outer sep=0pt,minimum width=0pt]\node[](#1){#2};}

\newcommand{\refines}[2]%
{\ensuremath{\textcolor{tgtcol}{#1} \slantsqsubseteq \textcolor{srccol}{#2}}}

% \setbeameroption{show only notes}
\setbeamerfont{note page}{size=\scriptsize}
\setbeamertemplate{footline}[frame number]{}
\setbeamertemplate{navigation symbols}{}
% \includeonlyframes{asymmetry}

\title{Relational Separation Logic for Compiler Verification}

\author{
  Gabriel Desfrene\\
  Supervised by Xavier Leroy and François Pottier
}

\date{September 2026}

\begin{document}

\begin{frame}[label=title]
  \titlepage

  \note[item]{
    Bonjour, je suis Gabriel Desfrene.
  }

  \note[item]{
    Je vais vous présenter le travail que j'ai réalisé sous la direction
    de Xavier Leroy et François Pottier, au sein de l'équipe Inria Cambium.
  }

  \note[item]{
    Ce travail porte sur le développement d'une logique de séparation relationnelle
    qui est dédiée à la vérification des compilateurs.
  }
\end{frame}

\section{Motivation}

\begin{frame}[label=bugs]
  \frametitle{Compilers Can Introduce Bugs}

  \begin{center}
    \begin{tikzpicture}[
        comp/.style={align=center, text width=3cm, font=\small},
        bug/.style={
          background default text=red!70!black,
          background text=gray!50,
        },
        >=stealth, shorten >=2pt
      ]
      \coordinate (topleft) at (-52mm, 20mm);
      \coordinate (bottomright) at (52mm, -40mm);
      \coordinate (topright) at (topleft -| bottomright);
      \coordinate (topmid) at ($(topleft)!0.5!(topright)$);
      \coordinate (bottommid) at (topmid |- bottomright);

      \draw[bbox] (topleft) rectangle (bottomright);

      \node<2->[draw, minimum height=1cm, minimum width=2.5cm, fill=gray!20]
      (compiler) at ($(topmid)!0.25!(bottommid)$) {Compiler};

      \node<2->[left=of compiler, text=srccol] (source) {Source Program};
      \node<2->[right=of compiler, text=tgtcol] (target) {Target Program};

      \draw<2->[->, thick] (source) -- (compiler);
      \draw<2->[->, thick]
      (compiler) -- node[above, bug, text on=<4->]
      {
        \begin{tikzpicture}[baseline]
          \node<3->[inner sep=0pt] (bug) at (0, 0) {\emoji{bug}};
          \only<4->{%
            \fill[white, fill opacity=0.6] (bug.south west) rectangle (bug.north east);
          }
        \end{tikzpicture}
      }
      (target);

      \node<4->[below=0.2cm of compiler, text=green!70!black]
      (verified) {\textbf{Verified}};

      \node<5->[comp] (compcert) at ($(topmid)!0.75!(bottommid)$)
      {CompCert\\ C $\to$ machine};

      \node<5->[left=1em of compcert, comp]
      (cakeml) {CakeML\\ ML $\to$ machine};

      \node<5->[right=1em of compcert, comp]
      (jasmin) {Jasmin\\ \small crypto $\to$ machine};
    \end{tikzpicture}
  \end{center}

  \note<1->[item]{
    Lorsque l'on execute un programme, ce n'est pas le code source
    qui est directement exécuté par la machine.
  }

  \note<2->[item]{
    On fait appel à un compilateur qui traduit ce programme source
    en un programme dans un langage intermédiaire de plus bas niveau,
    comprehensible par la machine.
    Pourquoi vérifier un compilateur ?
  }

  \note<3->[item]{
    Parce qu'un compilateur est lui-même un programme complexe
    qui peut introduire des bugs
    ou des comportements non désirés
    lors des transformations qu'il opère.
  }

  \note<4->[item]{
    Vérifier un compilateur consiste à vérifier
    que le programme effectuant ces transformations est correct.
  }

  \note<5->[item]{
    Il existe aujourd'hui plusieurs compilateurs vérifiés :
    CompCert qui compile du C vers un langage machine,
    CakeML qui compile un langage fonctionnel vers un langage machine,
    ou encore Jasmin, spécialisé dans les implémentations cryptographiques.
  }

\end{frame}


\begin{frame}[label=refinement]
  \frametitle{Behavioral Refinement}

  \begin{center}
    \begin{tikzpicture}[
      center box/.style={
          anchor=north,
          align=center, text width=10mm,
          background default text=black,
          background text=gray,
          % draw
        },
        src box/.style={
          align=center, text width=25mm,
          background default text=srccol,
          background text=gray,
          % draw
        },
        tgt box/.style={
          align=center, text width=25mm,
          background default text=tgtcol,
          background text=gray,
          % draw
        },
        beh box/.style={
          align=center, text width=6cm, font=\small, text=darkgray
        }
      ]
      \coordinate (topleft) at (-52mm, 25mm);
      \coordinate (bottomright) at (52mm, -45mm);
      \coordinate (topright) at (topleft -| bottomright);
      \coordinate (topmid) at ($(topleft)!0.5!(topright)$);
      \coordinate (bottommid) at (topmid |- bottomright);

      \draw[bbox] (topleft) rectangle (bottomright);

      \node<1->[center box, text on=<3->] (c1) at ($(topmid)!0.1!(bottommid)$)
      {$\implies$};
      \node<1->[left=1em of c1, src box, text on=<3->] (s1) {$S \vDash Spec$};
      \node<1->[right=1em of c1, tgt box, text on=<3->] (t1) {$T \vDash Spec$};

      \node<2->[anchor=south, beh box] (beh) at ($(topmid)!0.9!(bottommid)$)
      {
        $b \in \left\{
          \texttt{Term}(v, m) \mid \texttt{Div} \mid \text{\lightning}
        \right\}$
      };

      \node<3->[center box, text on=<4->] (c2) at ($(topmid)!0.2!(bottommid)$)
      {$\iff$};
      \node<3->[left=1em of c2, src box, text on=<4->] (s2) {$b \in S$};
      \node<3->[right=1em of c2, tgt box, text on=<4->] (t2) {$b \in T$};

      \node<4->[center box, text on=<5->] (c3) at ($(topmid)!0.3!(bottommid)$)
      {$\impliedby$};
      \node<4->[left=1em of c3, src box, text on=<5->] (s3) {$b \in S$};
      \node<4->[right=1em of c3, tgt box, text on=<5->] (t3) {$b \in T$};

      \coordinate<5-> (midup) at ($(topmid)!0.45!(bottommid)$);
      \coordinate<5-> (midlow) at ($(topmid)!0.75!(bottommid)$);

      \def\halfwidth{45mm}
      \coordinate<5-> (topleft) at ($ (midup.north) + (-\halfwidth, 0) $);

      \draw<5->[rounded corners=8pt, fill=gray!10, draw=gray!80, thick]
      (topleft) rectangle ($ (midlow.south) + (\halfwidth, 0) $);

      \node<5->[below right=1ex of topleft] (ref) {$\refines{T}{S} \triangleq {}$};

      \node<6->[center box, text on=<1>] (c4) at ($(topmid)!0.6!(bottommid)$)
      {$\impliedby$};
      \node<6->[left=1em of c4, src box, text on=<1>] (s4) {$b \in S \vee \text{\lightning} \in S$};
      \node<6->[right=1em of c4, tgt box, text on=<1>] (t4) {$b \in T$};
    \end{tikzpicture}
  \end{center}

  \note[item]{
    Comment formaliser la correction d'un compilateur ?
    On souhaite que toute garantie obtenue sur le programme source
    reste valable pour le programme produit.
  }

  \note[item]{
    Ces garanties portent sur les comportements observables du programme source.
    soit il termine avec une valeur et une mémoire finale,
    soit il diverge, soit il est bloqué dans un état indéfini (stuck).
  }

  \note[item]{
    Exiger l'égalité stricte des comportements est trop restrictif : un compilateur
    a le droit de restreindre le non-déterminisme, par exemple en fixant l'ordre
    d'évaluation des arguments. On garantit donc une inclusion des comportements.
  }

  \note[item]{
    De plus, face à un comportement indéfini (comme une division par zéro non utilisée),
    le compilateur peut éliminer l'opération.
  }

  \note[item]{
    La correction d'un compilateur s'énonce sous la forme d'un raffinement de
    comportement :
    tout comportement exhibé par le programme cible doit être un comportement
    possible du programme source, ou bien le source contenait un comportement indéfini.
  }

\end{frame}


\begin{frame}[label=simulations]
  \frametitle{Simulation Relations}

  \begin{center}
    \begin{tikzpicture}[
        title/.style={
          align=center, text width=20mm, minimum height=7mm,
          font=\bfseries,
          % draw
        },
        step arrow/.style={
          >=stealth, ->, thick
        },
        rel/.style={
          dashed, darkgray
        },
        conf/.style={
          align=center, text width=5mm, minimum height=6mm,
          % draw,
        },
        arr/.style={
          ->, >=stealth, thick, gray
        },
      ]

      \coordinate (topleft) at (-45mm, 20mm);
      \coordinate (bottomright) at (45mm, -40mm);
      \coordinate (topright) at (topleft -| bottomright);

      \draw[bbox] (topleft) rectangle (bottomright);

      \node<1> at ($(topleft)!0.5!(bottomright)$)
      {$\begin{array}{c}
        \textcolor{tgtcol}{t_0 \to t_1 \to \dots \to \mathtt{Term}(v, m)} \\
        \Downarrow \\
        \textcolor{srccol}{s_0 \to s_1 \to \dots \to \mathtt{Term}(v, m)} \\
        \end{array}
        $};

      \node<2->[title] (forward) at ($(topleft)!0.22!(topright) - (0, 2em)$)
      {\only<3->{Forward}};

      \node<2->[below=1em of forward.south west, conf, srccol] (sf1) {$s$};
      \node<2->[below=12mm of sf1, conf, srccol] (sf2) {$s'$};

      \draw<2->[step arrow, srccol]
      (sf1) -- node[left] (msf) {\only<4->{$\forall$}} (sf2);

      \node<2->[below=1em of forward.south east, conf, tgtcol] (tf1) {$t$};
      \node<2->[below=12mm of tf1, conf, tgtcol] (tf2) {$t'$};

      \draw<2->[step arrow, tgtcol]
      (tf1) -- node[right] (mtf) {\only<4->{$\exists$}} (tf2);

      \draw<2->[rel] (sf1) -- (tf1);
      \draw<2->[rel] (sf2) -- node[] (lowfor) {} (tf2);

      \draw<4->[arr] ($(msf) + (7mm, 0)$) to[bend left=25] ($(mtf) - (7mm, 0)$);

      \node<5->[below=1em of lowfor, font=\scriptsize] (dethyp) {${} + \text{Determinism}$};
      \node<5->[below=.5mm of dethyp] (arrfor) {$\Downarrow$};
      \node<5->[below=.5mm of arrfor] (reffor) {$\refines{T}{S}$};

      \node<6->[title] (backward) at ($(topleft)!0.78!(topright) - (0, 2em)$) {Backward};

      \node<6->[below=1em of backward.south west, conf, srccol] (sb1) {$s$};
      \node<6->[below=12mm of sb1, conf, srccol] (sb2) {$s'$};

      \draw<6->[step arrow, srccol]
      (sb1) -- node[left] (msb) {\only<7->{$\exists$}} (sb2);

      \node<6->[below=1em of backward.south east, conf, tgtcol] (tb1) {$t$};
      \node<6->[below=12mm of tb1, conf, tgtcol] (tb2) {$t'$};

      \draw<6->[step arrow, tgtcol] (tb1) -- node[right]
      (mtb) {\only<7->{$\forall$}} (tb2);

      \draw<6->[rel] (sb1) -- (tb1);
      \draw<6->[rel] (sb2) -- node[] (lowback) {}(tb2);

      \draw<7->[arr] ($(mtb) - (7mm, 0)$) to[bend right=25] ($(msb) + (7mm, 0)$);

      \node<8-> (refback) at (reffor -| lowback) {$\refines{T}{S}$};
      \node<8-> (arrback) at ($(lowback)!0.5!(refback)$) {$\Downarrow$};

    \end{tikzpicture}
  \end{center}

  \note<1->[item]{
    Ce raffinement est énoncé sur des traces d'exécution complètes.
    Mais prouver directement une propriété sur des exécutions globales est difficile,
    car les optimisations opérent localement, sur quelques étapes de calcul.
  }

  \note<2->[item]{
    On utilise alors les relations de simulation.
    Elles relient pas à pas les configurations successives des deux programmes.
  }

  \note<3->[item]{
    Les simulations dites 'en avant' relient chaque pas du source à des pas
    de la cible.
  }

  \note<5->[item]{
    Elles garantissent le raffinement lorsque le langage cible est déterministe.
  }

  \note<6->[item]{
    Une autre possibilité sont les simulations dites 'en arrière'.
  }

  \note<7->[item]{
    On observe les pas du programme cible et on les relie à des pas du source.
  }

  \note<8->[item]{
    Une simulation en arrière implique directement le raffinement sans hypothèse
    de déterminisme sur le langage cible.
  }

\end{frame}


\begin{frame}[label=stuttering]
  \frametitle{The Stuttering Problem}

  \begin{center}
    \begin{tikzpicture}[
        step arrow/.style={
          >=stealth, ->, thick
        },
        rel/.style={
          dotted, darkgray
        },
        conf/.style={
          align=center, text width=5mm, minimum height=6mm,
          % draw,
        },
        arr/.style={
          ->, >=stealth, thick, gray
        },
        fuel/.style={
          darkgray, font=\small
        },
      ]
      \coordinate (topleft) at (-45mm, 30mm);
      \coordinate (bottomright) at (45mm, -40mm);
      \coordinate (topright) at (topleft -| bottomright);
      \coordinate (topmid) at ($(topleft)!0.5!(topright)$);

      \draw[bbox] (topleft) rectangle (bottomright);

      \node[conf, srccol] (s0) at ($(topmid) + (-10mm, -10mm)$) {$s_0$};
      \node[below=10mm of s0, conf, srccol] (s1) {};
      \node[below=10mm of s1, conf, srccol] (s) {$s$};

      \node[conf, tgtcol] (t0) at ($(topmid) + (10mm, -10mm)$) {$t_0$};
      \node[below=10mm of t0, conf, tgtcol] (t1)
      {\only<2->{$t_1$}};
      \node[below=10mm of t1, conf, tgtcol] (t)
      {\alt<5->{$t_2$}{$t$}};

      % Frame 1
      \draw<1->[step arrow, srccol] (s0) -- (s);
      \draw<1->[rel] (s0) -- (t0);
      \draw<1->[rel] (s) -- (t);

      \draw<1>[step arrow, tgtcol] (t0) -- (t);

      % Frame 2
      \draw<2->[step arrow, dashed, tgtcol] (t0) -- node (stut) {} (t1);
      \draw<2->[step arrow, tgtcol] (t1) --  (t);

      % Frame 3
      \draw<3-4>[rel] (s1) -- node[above, gray] (qmark) {?} (t1);

      % Frame 4
      \node<4> (stutlbl) at (stut.east -| {$(topleft)!0.85!(topright)$}) {Stuttering};
      \draw<4>[gray, ->] (stutlbl) -- (stut.east);

      % Frame 5
      \node<5->[below=5mm of t, conf, tgtcol] (t3) {$\dots$};
      \node<5->[below=5mm of t3, conf, tgtcol] (tinf)
      {\alt<7->{$t$}{$\dots$}};

      \draw<5->[step arrow, tgtcol, dashed] (t) -- (t3);
      \draw<5->[step arrow, tgtcol, dashed] (t3) -- (tinf);

      % Frame 6
      \node<6->[below=5mm of s, conf, srccol] (s3) {};
      \node<6->[fuel] (f1) at ($(s0)!0.5!(t1)$) {$1$};
      \node<6->[fuel] (f2) at ($(s1)!0.5!(t)$) {$0$};
      \node<6->[fuel] (le1) at ($(f1)!0.5!(f2)$) {$\vee$};

      \node<6->[fuel] (f3) at ($(s)!0.5!(t3)$) {$n$};
      \node<6->[fuel] (f4) at ($(s3)!0.5!(tinf)$) {$n-1$};
      \node<6->[fuel] (le2) at ($(f3)!0.5!(f4)$) {$\vee$};

      % Frame 7
      \node<7->[below=5mm of s3, conf, srccol] (sinf) {$s$};
      \draw<7->[thick, double equal sign distance, srccol] (s) -- (sinf);

      \draw<7->[rel] (sinf) -- (tinf);
    \end{tikzpicture}
  \end{center}

  \note<1->[item]{
    Une simulation ne peut pas simplement associer un pas de calcul cible
    à exactement un pas de calcul source.
  }

  \note<2->[item]{
    En effet, un programme source dans un langage de plus haut peut faire moins d'etapes
    que sont transformé dans un language de plus bas niveau.
    Contrairement, un programme cible optimisé peut réalisé moins d'étapes que
    son programme source associé.
  }

  \note<3->[item]{
    Lorsque l'un dees deux programme réalise un pas de calcul qui n'a pas
    d'équivalent dans l'autre programme,
  }

  \note<4->[item]{
    on appèle cela une étape de bégaiement.
  }

  \note<5->[item]{
    Cependant, ce bégaiement doit être strictement contrôlé :
    si un nombre infini d'étapes de bégaiement est autorisé,
    un programme source divergent peut être relié
    à n'importe quelle cible.
    Cela n'est pas permis par la relation de raffinement.
  }

  \note<6->[item]{
    Une solution classique est de fixer une mesure bien fondée qui décroît strictement
    à chaque pas de bégaiement et se réinitialise lors d'un pas commun.
  }

  \note<7->[item]{
    Cela assure que le nombre d'étapes de bégaiement
    entre chaque synchronisation est borné.
  }

  \note<6->[item]{
    Cependant, construire explicitement cette mesure pour chaque transformation
    peut se réveler contraignant voire difficile.
  }

\end{frame}


\begin{frame}[label=limitations]
  \frametitle{Limitations of Existing Approaches}

  \begin{center}
    \begin{tikzpicture}[
        header/.style={
          anchor=north, text=blue!85!black
        },
        side/.style={
          anchor=west, font=\bfseries\small, text=gray!80!black
        },
        elem-bad/.style={
          font=\small, text=red!75!black,
        },
        elem-warn/.style={
          font=\small, text=orange!80!black,
        },
        missing/.style={
          font=\bfseries\small,
          minimum height=1cm, text width=30mm, align=center,
          fill=green!15, draw=green!60!black, text=green!60!black,
          thick, rounded corners=6pt, inner sep=5pt,
        },
        gridline/.style={
          draw=gray!30, thick
        }
      ]
      \coordinate (topleft) at (-52mm, 25mm);
      \coordinate (bottomright) at (52mm, -25mm);
      \coordinate (topright) at (topleft -| bottomright);
      \coordinate (bottomleft) at (topleft |- bottomright);

      \coordinate (col1) at ($(topleft)!0.22!(topright)$);
      \coordinate (col2) at ($(topleft)!0.61!(topright)$);

      \draw[bbox] (topleft) rectangle (bottomright);

      \draw<2->[gridline] (col1) -- (col1 |- bottomright);
      \draw<2->[gridline] (col2) -- (col2 |- bottomright);

      \node<3->[header, minimum width=36mm] (nd) at ($(col1)!0.5!(col2)$) {Non-Determinism};
      \coordinate<3-> (l1) at (nd.south -| topleft);
      \draw<3->[gridline] (l1) -- (l1 -| topright);

      \node<4->[side] (cakeml) at ($(l1)!0.25!(bottomleft)$) {CakeML};
      \node<4->[elem-bad] at (cakeml -| nd) {Not Supported};

      \node<5->[side] (compcert) at ($(l1)!0.75!(bottomleft)$) {CompCert};

      \coordinate<5-> (l2) at ({$(cakeml)!0.5!(compcert)$} -| l1);
      \draw<5->[gridline] (l2) -- (l2 -| topright);
      \node<5->[elem-warn] at (compcert -| nd) {External only};

      \node<6->[missing] at (l2 -| nd) {Backward Simulation};

      \node<7->[header, minimum width=36mm] (mem) at ($(col2)!0.5!(topright)$) {Memory Reasoning};

      \node<8->[elem-warn] at (cakeml -| mem) {Global Invariants};
      \node<8->[elem-warn] at (compcert -| mem) {Global Invariants};

      \node<9->[missing] at (l2 -| mem) {Separation Logic};

    \end{tikzpicture}
  \end{center}

  \note<2->[item]{
    Les approches existantes présentent deux limitations principales.
  }

  \note<3->[item]{
    La première est la gestion du non-déterminisme interne.
  }

  \note<4->[item]{
    CakeML exige un langage cible déterministe.
  }

  \note<5->[item]{
    CompCert ne traite que le non-déterminisme externe.
  }

  \note<6->[item]{
    En effet, ils s'appuient en partie sur des simulations en avant.
    Traiter de la concurrence ou des effets
    non déterministes demande une simulation en arrière.
  }

  \note<7->[item]{
    La deuxieme est le raisonnement sur la mémoire.
  }

  \note<8->[item]{
    CompCert et CakeML manipulent tous les deux des invariants globaux sur la mémoire.
    Cela complique les preuves dès que l'organisation du tas change, meme un petit peu.
  }

  \note<9->[item]{
    La logique de séparation permet de raisonner localement sur des segments
    de mémoire disjoints,
    mais son intégration dans des compilateurs vérifiés reste encore à développer.
  }

\end{frame}

\begin{frame}[label=contributions]
  \frametitle{Our Contributions}

  \begin{center}
    \begin{tikzpicture}[
        layerbox/.style={
          align=center, minimum height=10mm,
          thick, rounded corners=6pt, inner sep=4pt
        },
        freebox/.style={
          layerbox, minimum width=70mm,
          fill=green!10, draw=green!60!black, text=green!50!black,
        },
        logicbox/.style={
          layerbox, minimum width=70mm,
          fill=teal!12, draw=teal!70!black, text=teal!75!black,
        },
        langbox/.style={
          layerbox, minimum width=70mm,
          fill=blue!10, draw=blue!70!black, text=blue!80!black,
        },
        transbox/.style={
          layerbox,
          fill=orange!10, draw=orange!70!black, text=orange!70!black,
        },
        optimbox/.style={
          layerbox,
          fill=violet!10, draw=violet!70!black, text=violet!70!black,
        },
        scope-tag/.style={
          font=\tiny\itshape, align=left, text=gray!70!black, inner xsep=0pt
        },
        divider/.style={
          draw=gray!40, thick, dashed
        }
      ]

      \coordinate (topleft) at (-52mm, 25mm);
      \coordinate (bottomright) at (52mm, -30mm);
      \coordinate (topright) at (topleft -| bottomright);
      \coordinate (bottomleft) at (topleft |- bottomright);
      \coordinate (topmid) at ($(topleft)!0.5!(topright)$);
      \coordinate (bottommid) at (topmid |- bottomright);

      \draw[bbox] (topleft) rectangle (bottomright);

      \node<2->[freebox, anchor=south] (free) at (bottommid)
      {Free Backward Simulation};

      \node<3->[above=1.5mm of free, logicbox] (rsl)
      {Relational Separation Logic};

      \coordinate<3->[yshift=1.5mm] (sep) at (rsl.north);
      \coordinate<3-> (sepleft) at (topleft |- sep);
      \coordinate<3-> (sepright) at (bottomright |- sep);

      \draw<4->[divider] (sepleft) -- (sepright);
      \node<4->[anchor=north west, scope-tag] at (sepleft)
      {Language\\ Independent};
      \node<4->[anchor=south west, scope-tag] at (sepleft)
      {Language\\ Dependent};

      \node<5->[above=3mm of rsl, langbox] (murtl)
      {Low level language: \muRTL{}};

      \coordinate<6-> (l) at (murtl.north west);
      \coordinate<6-> (r) at (murtl.north east);
      \coordinate<6-> (m1) at ($(l)!0.66!(r)$);
      \coordinate<6->[xshift=1.5mm] (m2) at (m1);

      \node<6->[
        above=1.5mm of murtl.north west, anchor=south west,
        transbox, text width between={l and m1}
      ]
      (transf) {Program Transformations};

      \node<7->[
        above=1.5mm of murtl.north east, anchor=south east,
        optimbox, text width between={m2 and r}
      ]
      (optim) {Tunneling};

    \end{tikzpicture}
  \end{center}

  \note<1->[item]{
    Pour répondre à ces limitations, nous avons développé une infrastructure
    pour la vérification de compilateurs.
  }

  \note<2->[item]{
    Nous avons adapté une relation de simulation en arrière dite 'libre'.
    Celle-ci ne demande aucune mesure explicite pour borner le bégaiement
    et évite de coupler des étapes de calcul.
  }

  \note<3->[item]{
    Nous avons développé une logique de séparation relationnelle
    pour raisonner localement sur des paires de mémoires disjointes.
  }

  \note<5->[item]{
    Nous avons instancié ce cadre génerique pour un langage intermédiaire de bas niveau,
    \muRTL{}.
  }

  \note<6->[item]{
    Puis nous avons vérifiés des transformations de programmes.
  }

  \note<7->[item]{
    Enfin, nous avons vérifiés une transformation globale
    de CompCert: le raccourcissement des sauts.
    L'ensemble a été formalisé et vérifié en Rocq.
  }

\end{frame}


\section{A Simulation}

\begin{frame}[label=free]
  \frametitle{The Free Simulation}

  \begin{center}
    \begin{tikzpicture}[
        arr/.style={
          >=stealth, ->, thick,
          background aspect=densely dash dot,
        },
        step arrow/.style={
          arr, background default aspect=solid,
        },
        stut arrow/.style={
          arr, background default aspect=dashed,
        },
        rel/.style={
          dotted, darkgray
        },
        conf/.style={
          anchor=north,
          align=center, text width=5mm, minimum height=6mm,
          % draw,
        },
        legend/.style={
          gray, font=\scriptsize,
          draw=none, fill=white, text=gray,
        },
        legend arr/.style={
          ->, >=stealth, thick, gray
        },
        fuel/.style={
          anchor=center, gray, font=\footnotesize,
          % draw
        },
        idx legend/.style={
          legend, anchor=west, text width=8mm, align=center
        }
      ]

      \coordinate (topleft) at (-45mm, 35mm);
      \coordinate (bottomright) at (45mm, -35mm);
      \coordinate (topright) at (topleft -| bottomright);
      \coordinate (topmid) at ($(topleft)!0.5!(topright)$);

      \def\halfwidth{20mm}
      \def\stepsize{15mm}

      \draw[bbox] (topleft) rectangle (bottomright);

      % Frame 2
      \node<2->[conf, srccol, xshift=-\halfwidth] (s0) at (topmid) {$s_0$};
      \node<2->[below=2*\stepsize of s0.north, conf, srccol] (s1) {$s_1$};
      \node<2->[below=\stepsize of s1.north, conf, srccol] (s2) {$s_2$};
      \node<2->[below=\stepsize of s2.north, conf, srccol] (s3) {$s_3$};

      \node<2->[conf, tgtcol, xshift=\halfwidth] (t0) at (topmid) {$t_0$};
      \node<2->[below=\stepsize of t0.north, conf, tgtcol] (t1) {$t_1$};
      \node<2->[below=\stepsize of t1.north, conf, tgtcol] (t2) {$t_2$};
      \node<2->[below=2*\stepsize of t2.north, conf, tgtcol] (t3) {$t_3$};

      \draw<2->[step arrow, srccol, aspect on=<7->] (s0) -- node (syns1) {} (s1);
      \draw<2->[stut arrow, srccol, aspect on=<7->] (s1) -- node (begs) {} (s2);
      \draw<2->[step arrow, srccol, aspect on=<7->] (s2) -- node (syns2) {} (s3);

      \draw<2->[stut arrow, tgtcol, aspect on=<7->] (t0) -- node (begt) {} (t1);
      \draw<2->[step arrow, tgtcol, aspect on=<7->] (t1) -- node (synt1) {} (t2);
      \draw<2->[step arrow, tgtcol, aspect on=<7->] (t2) -- node (synt2) {} (t3);

      % Frame 3
      \node<3-5>[legend] (begt-leg) at ({$(topleft)!0.9!(topright)$} |- begt) {Stuttering};
      \draw<3-5>[legend arr] (begt-leg) -- (begt);

      \node<3-5>[legend] (begs-leg) at ({$(topleft)!0.1!(topright)$} |- begs) {Stuttering};
      \draw<3-5>[legend arr] (begs-leg) -- (begs);

      % Frame 4
      \node<4-5>[legend] (syn1-leg) at ($(syns1)!0.5!(synt1)$) {Synchronous Step};
      \draw<4-5>[legend arr] (syn1-leg) -- (syns1);
      \draw<4-5>[legend arr] (syn1-leg) -- (synt1);

      \node<4-5>[legend] (syn2-leg) at ($(syns2)!0.5!(synt2)$) {Synchronous Step};
      \draw<4-5>[legend arr] (syn2-leg) -- (syns2);
      \draw<4-5>[legend arr] (syn2-leg) -- (synt2);

      % Frame 5
      \draw<5, 8->[rel] (s0) -- node [legend] {Related} (t0);
      \draw<5, 8->[rel] (s1) -- node [legend] {Related} (t2);
      \draw<5, 8->[rel] (s3) -- node [legend] {Related} (t3);

      % Frame 6
      % Clear the labels

      % Frame 7
      % The steps change to dash dot

      % Frame 8
      \node<8->[fuel, left=0mm of syns1] {1};
      \node<8->[fuel, left=0mm of begs] {1};
      \node<8->[fuel, left=0mm of syns2] {2};

      \node<8->[fuel, right=0mm of begt] {1};
      \node<8->[fuel, right=0mm of synt1] {2};
      \node<8->[fuel, right=0mm of synt2] {1};

      % Frame 9
      \node<9->[idx legend] (idx) at (s1 -| topleft) {Index\\ Reset};
      \draw<9->[legend arr] (idx) to[out=90, in=180] (s0);
      \draw<9->[legend arr] (idx) -- (s1);
      \draw<9->[legend arr] (idx) to[out=-90, in=180] (s3);

    \end{tikzpicture}
  \end{center}

  \note[item]{
    Notre relation de simulation est adaptée des travaux de Cho et Song:
    \emph{Stuttering for Free}.
  }

  \note<2->[item]{
    Dans les simulations traditionnelles,
  }

  \note<3->[item]{
    les étapes de bégaiement doivent être suivies
  }

  \note<4->[item]{
    d'une étape commune
  }

  \note<5->[item]{
    pour permettre de relier les configuration poursuivre de la preuve.
  }

  \note<6->[item]{
    La simulation libre dévelopée dans \emph{Stuttering For Free}
    oblige une progression entièrement asymétrique.
  }

  \note<7->[item]{
    Les deux programmes progressent de manière indépendante.
  }

  \note<8->[item]{
    Pour empêcher le bégaiement infini,
    la simulation enregistre via des indices le nombre de pas effectués par
    chaque programme depuis les derniers états reliés.
  }

  \note<9->[item]{
    Dès que les deux programmes ont chacun progressé au moins une fois,
    on peut relier les états.
    Les indices sont alors réinitialisées et on peut poursuivre la preuve.
  }

  \note<9->[item]{
    Intuitivement, au lieu de borner à l'avance le nombre de pas
    de bégaiement restants, on comptabilise le progrès déjà réalisé.
  }

  \note<9->[item]{
    Ce découplage entre synchronisation et progression simplifie
    l'écriture des invariants et le raisonnement pas à pas.
  }

\end{frame}


\def\lowcol{darkgray}
\newcommand{\fsimrel}[4]{
  \textcolor{\lowcol}{\langle} \textcolor{tgtcol}{#1}\textcolor{\lowcol}{, #2\rangle}
  \sim_{\text{free}}
  \textcolor{\lowcol}{\langle} \textcolor{srccol}{#4}\textcolor{\lowcol}{, #3 \rangle}
}

\begin{frame}[label=theory]
  \frametitle{Theoretic Guarantees}

  \pause

  \textcolor{gray}{Sound:}
  \begin{equation*}
    \textcolor{\lowcol}{\forall i\, j,\qquad }
    \fsimrel{T}{i}{j}{S} \implies \refines{T}{S}
  \end{equation*}

  \vspace{5mm}
  \pause

  \textcolor{gray}{Equivalent:}
  \begin{equation*}
    (\textcolor{\lowcol}{\exists i\, j,\ } \fsimrel{T}{i}{j}{S})
    \iff \textcolor{tgtcol}{T} \sim_{\text{explicit}} \textcolor{srccol}{S}
  \end{equation*}

  \vspace{5mm}
  \pause

  \def\lowcol{black}

  \textcolor{gray}{Index Irrelevant:}
  \begin{equation*}
    (\textcolor{\lowcol}{\exists i\, j,\ } \fsimrel{T}{i}{j}{S})
    \implies
    \textcolor{\lowcol}{\forall i\, j,\quad } \fsimrel{T}{i}{j}{S}
  \end{equation*}

  \note<1->[item]{
    La flexibilité apportée par cette simulation n'est pas au détriment
    de son expressivité.
  }

  \note<2->[item]{
    Cette simulation implique le raffinement de comportement souhaité.
  }

  \note<3->[item]{
    De plus, elle est équivalente aux simulations explicites avec mesures bien fondées.
  }

  \note<4->[item]{
    Un corrolaire de cette équivalence montre également que les indices de départ
    n'ont aucune incidence sur la validité de la simulation.
  }

  \note<4->[item]{
    Ils ne servent qu'à suivre le progrès de chaque programme lors de la preuve.
  }

\end{frame}

\def\treesibdist{12mm}
\def\treelevdist{12mm}

\begin{frame}[label=asymmetry]
  \frametitle{Asymmetric Reasoning}

  \begin{center}
    \begin{tikzpicture}[
        header/.style={
          anchor=north, minimum height=8mm, minimum width=40mm,
          % draw
        },
        gridline/.style={
          draw=gray!30, thick
        },
        link/.style={
          thick, ->, >=stealth,
        }
      ]
      \coordinate (topleft) at (-52mm, 30mm);
      \coordinate (bottomright) at (52mm, -30mm);
      \coordinate (topright) at (topleft -| bottomright);
      \coordinate (bottomleft) at (topleft |- bottomright);
      \coordinate (topmid) at ($(topleft)!0.5!(topright)$);
      \coordinate (bottommid) at (topmid |- bottomright);

      \draw[bbox] (topleft) rectangle (bottomright);

      \node<2->[header, tgtcol] (tgt-head) at ($(topleft)!0.5!(topmid)$) {Target};
      \node<2->[header, srccol] (src-head) at ($(topmid)!0.5!(topright)$) {Source};

      \coordinate<2-> (l1) at (tgt-head.south -| topleft);
      \coordinate<2->[yshift=15mm] (midend) at (topmid |- bottomright);

      \draw<2->[gridline] (l1) -- (l1 -| topright);
      \draw<2->[gridline] (topmid) -- (midend);

      \node<3->[below=1mm of tgt-head, font=\small] (tgt-nd)
      {\emoji{smiling-face-with-horns} nondeterminism};

      % Target tree
      \node<4->[below=1mm of tgt-nd, tgtcol] (roott) {$\forall$};
      \node<4-> (t-l) at ($(roott) + (-0.5*\treesibdist, -\treelevdist)$) {};
      \node<4->[tgtcol] (t-r) at ($(roott) + (0.5*\treesibdist, -\treelevdist)$)
      {$\forall$};
      \node<4-> (endt) at ($(t-r) + (-0.5*\treesibdist, -\treelevdist)$) {};
      \node<4-> (t-rr) at ($(t-r) + (0.5*\treesibdist, -\treelevdist)$) {};

      \draw<4->[link, tgtcol] (roott) -- (t-l);
      \draw<4->[link, tgtcol] (roott) -- (t-r);
      \draw<4->[link, tgtcol] (t-r) -- (endt);
      \draw<4->[link, tgtcol] (t-r) -- (t-rr);

      \node<5->[below=1mm of src-head, font=\small] (src-nd)
      {\emoji{smiling-face-with-halo} nondeterminism};

      % Source Tree
      \node<6->[below=1mm of src-nd, srccol] (roots) {$\exists$};
      \node<6-> (s-l) at ($(roots) + (-0.5*\treesibdist, -\treelevdist)$) {};
      \node<6->[srccol] (s-r) at ($(roots) + (0.5*\treesibdist, -\treelevdist)$)
      {$\exists$};
      \node<6-> (ends) at ($(s-r) + (-0.5*\treesibdist, -\treelevdist)$) {};
      \node<6-> (s-rr) at ($(s-r) + (0.5*\treesibdist, -\treelevdist)$) {};

      \draw<6->[link, gray]   (roots) -- (s-l);
      \draw<6->[link, srccol] (roots) -- (s-r);
      \draw<6->[link, srccol] (s-r)   -- (ends);
      \draw<6->[link, gray]   (s-r)   -- (s-rr);


      \node<7-> at ($(midend)!0.5!(bottommid)$) {
        $\text{\lightning} \in \textcolor{srccol}{S}
        \implies
        \textcolor{tgtcol}{T} \sim_{\text{free}} \textcolor{srccol}{S}
        $
      };
    \end{tikzpicture}
  \end{center}

  \note<1->[item]{
    Puisque le raffinement exige que les comportements du code cible soient inclus
    dans ceux du source,
  }

  \note<2->[item]{
    les deux programmes sont traités de manière asymétrique.
  }

  \note<3->[item]{
    Le programme cible est traité avec un non-déterminisme démonique :
  }

  \note<4->[item]{
    On doit prouver qu'il peut progresser et établir la simulation
    pour toutes les transitions possibles.
  }

  \note[item]{
    Le programme source est quant à lui traité avec un non-déterminisme angélique :
    on choisit la branche d'exécution source appropriée pour correspondre à la cible.
  }

  \note[item]{
    Enfin, lorsque le programme source est bloqué,
    la simulation est immédiatement validée.
    Le compilateur est libre pour le code cible.
  }

\end{frame}

\section{A Logic}

\begin{frame}[label=logic]{Relational Separation Logic}
  \note[item]{
    La logique de séparation permet d'exprimer
    des propriétés locales sur un tas unique.
  }

  \note[item]{
    Dans notre cadre, nous devons raisonner simultanément sur deux tas :
    le tas du programme source et celui du programme cible.
  }

  \note[item]{
    Notre logique de séparation relationnelle manipule des prédicats
    sur ces deux mémoires.
  }

  \note[item]{
    La conjonction séparante relationnelle $P \sep Q$ décompose simultanément
    les deux tas en sous-tas disjoints,
    l'un sur lequel $P$ est valide et l'autre où $Q$ est valide.
  }

  \note[item]{
    On se donne deux assertions primitives décrivent les cellules allouées
    et libérées, tant côté source que côté cible.
  }

  \note[item]{
    Notre logique est linéaire.
    On ne peut pas ignorer une partie du tas dans un raisonnement,
    comme dans les logique d'Iris.
  }

  \note[item]{
    Enfin, notre logique est compatible avec le Proof Mode de Iris,
    ce qui facilite l'écriture des preuves.
  }

\end{frame}


\begin{frame}[label=values]{Values equivalence}
  \note[item]{
    Les transformations opérées par un compilateur
    changent souvent l'agencement de la mémoire
    et les adresses allouées.
  }

  \note[item]{
    Pour relier les deux espaces mémoire,
    nous introduisons une bijection partielle $I$
    entre les adresses cibles et les adresses sources.
  }

  \note[item]{
    Cette bijection est monotone.
    Elle grandie au fil de l'executiuon des programmes et des allocations.
    Elle sert à définir l'équivalence entre valeurs cibles et sources.
  }

  \note[item]{
    Deux valeurs sont équivalentes si :
    ce sont des entiers ou booléens égaux ;
    ce sont des pointeurs reliés par la bijection $I$ ;
    ou bien la valeur source est indéfinie ($\mathtt{undef}$),
    auquel cas la valeur cible peut être arbitraire.
  }

\end{frame}


\begin{frame}[label=injection]{Memory Injection}
  \note[item]{
    La bijection $I$ permet de définir le prédicat d'injection de mémoire $\meminj{I}{E}$,
    qui formalise la cohérence globale entre les deux tas. % TODO: Cite Simuliris.
  }

  \note[item]{
    L'assertion assure que pour chaque paire d'adresses dans $I$,
    les deux cellules sont soit allouées avec des valeurs équivalentes,
    soit toutes deux libérées.
  }

  \note[item]{
    Une règle d'extension permet d'ajouter de nouvelles adresses dans $I$,
    si elle respectent les condition précedentes.
  }

  \note[item]{
    Cependant,
    pour justifier un accès mémoire local dans les programmes,
    nous devons manipuler les assertions primitives correspondantes.
  }

  \note[item]{
    Retirer directement une paire d'adresses de la bijection $I$
    n'est pas possible :
    réduire $I$ pourrait invalider l'équivalence d'autres pointeurs du tas.
  }

  \note[item]{
    C'est le rôle de l'ensemble d'exceptions $E$.
    Deux lemmes d'extraction et de réintégration permettent de sortir temporairement
    une paire d'addresse de l'injection, puis de la réintégrer après modification locale.
  }


\end{frame}


\begin{frame}[label=hoare]{Relational Hoare Quadruples}
  \note[item]{
    Nous définissons des quadruplets de Hoare relationnels
    pour spécifier le raffinement entre paires de fonctions.
  }

  \note[item]{
    Leur sémantique est la suivantes:
    Si les arguments et mémoires initiales satisfont la précondition $P$
    et que le programme source n'a pas de comportement indéfinit:
    alors si la cible peut terminer,
    la source doit également terminer avec la postcondition $Q$.
    Si la cible peut diverger, la source doit également diverger.
  }

  \note[item]{
    Ce jugement est persistant dans notre logique de séparation :
    une fois prouvé pour une paire de fonctions,
    il peut être réutilisé autant de fois que nécessaire.
  }

\end{frame}

\section{A Language}

\begin{frame}[label=murtl]{\muRTL{} --- A Low-Level Intermediate Language}

  \note[item]{
    Pour tester ces outils, nous avons utilisé un langage
    intermédiaire inspiré de RTL (Register Transfer Language) dans CompCert:
    \muRTL{}.
  }

  \note[item]{
    Ce langage est non structuré : le corps d'une fonction est représenté par un graphe
    de flot de contrôle (CFG), où chaque instruction possède un label (PC)
    et indique ses successeurs directs.
  }

  \note[item]{
    Les instructions manipulent des registres, qui agissent comme des variables locales.
    Elles effectuent des opérations usuelles :
    des operations arithmétiques,
    des accès mémoires,
    des branchements conditionnels.
  }

  \note[item]{
    Chaque fonction possède ses propres registres locaux.
    Lors d'un appel de fonction,
    les registres de l'appelant sont sauvegardés dans la pile d'appel,
    au sein d'un cadre d'activation.
  }

\end{frame}


\begin{frame}[label=semantics]{Operational Semantics}

  \note[item]{
    La sémantique opérationnelle de \muRTL{} est exprimée à petits pas.
    Cette sémantique, décris comment la con figuration initiale évolue.
  }

  \note[item]{
    Une configuration est un triplet composé de la pile d'appels courante,
    de la mémoire et de l'état d'exécution du programme.
    Trois cas sont possibles pour l'états.
  }

  \note[item]{
    Au début du programme ainsi qu'apres un appel de fonction,
    le programe est dans l'état $\mathtt{CallState}$.
    Celui-ci indique la prochaine fonction à executer ainsi que la valeurs des arguments.
  }

  \note[item]{
    Au sein de la fonction le programme est dans un état classique $\mathtt{State}$
    qui indique la fonction courante, la prochaine instruction à executer ainsi que
    la valeurs des registres locaux.
  }

  \note[item]{
    A la fin de l'execution de la fonction, le programme passe dans
    l'état $\mathtt{ReturnState}$ retourne une valeur.
  }

  \note[item]{
    Cette sémantique permet de raisoner de manière modulaire sur la pile d'appel.
  }

  \note[item]{
    Pour prouver un quadruplet de hoare entre deux fonction,
    il suffit de le montrer dans une pile vide.
    Lors de l'utilisation de ce quadruplet,
    la pile d'appel des deux fonctions est augmentée par la pile courante.
  }

\end{frame}

\section{Case Studies}

\begin{frame}[label=zerocheck]{Case Study: Eliminating a Redundant Zero-Check}

  \note[item]{
    Avec ces outils, nous avons pu verifier des transformations de programmes concrètes.
    Une première transformation simple est l’élimination d’un test à zéro redondant.
  }

  \note[item]{
    Considérons une séquence source où une division est immédiatement suivie
    d'un test vérifiant si le diviseur $d$ est nul.
  }

  \note[item]{
    En \muRTL{}, la division par zéro est un état bloquant.
    Donc, lorsque le programme source va executer le test à zero,
    $d$ est obligatoirement non nul.
  }

  \note[item]{
    On peut donc optimizer ce code en remplaçant directement le test par
    une affectation par $\mathtt{false}$.
  }

  \note[item]{
    La preuve exploite l'asymétrie entre la source et la cible.
  }

  \note[item]{
    si $d = 0$, la source bloque et la cible a toute liberté.
  }

  \note[item]{
    si $d \neq 0$, la division réussie, et le test source vaut toujours faux.
    les deux programmes affecte au registre $t$ la meme valeur.
  }
\end{frame}


\begin{frame}[label=hoist]{Case Study: Hoisting an Invariant Load}
  \note[item]{
    Notre deuxième étude de cas porte sur deux fonction cette fois-ci.
  }

  \note[item]{
    On considère le programme ($\mathtt{fact\_slow}$) suivant.
    Celui-ci calcule la factorielle de $n$ dans $res$.
    Le décrement de $n$ est fait en lisant en memoire à l'adresse $addr$
    la valeur $1$ à chaque tour de boucle.
  }

  \note[item]{
    Dans le programme transformé $\mathtt{fact\_fast}$),
    ce chargement est optimisé dans le prélude de la fonction,
    et la boucle réutilise la valeur stockée dans un registre.
  }

  \note[item]{
    Cette transformation est correcte car la cellule mémoire pointée par $addr$
    n'est jamais modifiée pendant la boucle :
    la lire une unique fois au début est donc suffisant.
  }

\end{frame}


\begin{frame}[label=hoistproof]{?}
  \note[item]{
    La correction de cettre transformation est exprimée sous la forme
    d'un quadruplet de hoare.
  }

  \note[item]{
    Celui-ci demande en précondition que les arguments des fonctions soient équivalents,
    pour une bijection $I$ fixée. On a également le fait que la mémoire source et cible
    sont reliées selon $I$.
  }

  \note[item]{
    Si les programmes termine, on assure que les resultats sont équivalents pour la même
    bijection $I$. Les mémoires source et cible sont encore reliées selon $I$.
  }

  \note[item]{
    La preuve de cette optimisation se fait ici en maintenant l'invariant que
    la cellule pointée par $addr$ est toujours 1 lors de l'execution de la boucle.
  }
\end{frame}


\begin{frame}[label=tunneling]{Case Study: Short-Circuiting Jump Chains}
  \note[item]{
    Enfin, nous avons vérifié une transformation génerique de CompCert :
    le raccourcissement des chaînes de sauts (nop tunneling).
  }

  \note[item]{
    L'optimisation inspecte le graphe de contrôle et fait pointer chaque instruction
    directement vers son prochain successeur utile,
    en court-circuitant les suites de saut unitiles.
  }

  \note[item]{
    Il s'agit d'une transformation qui réécrit l'intégralité du graphe
    de flot de controle.
  }

  % Little more info on how this is defined ?

  \note[item]{
    La correction de cette passe d'otimisation est enoncee par le quadruplet
    de hoare suivant.
  }

  \note[item]{
    La preuve est faite pour une fonction source arbitraire.
    Elle utilise un invariant sur l'états des configuration ainsi qu'une
    une reccurence sur la profondeur à laquelle l'optimisation est faite.
  }

  \note[item]{
    Lorsque les deux programmes appellent une même fonction,
    l'invariant utilisé lors de la preuve est réutilisée directement,
    de manière modulaire.
  }

\end{frame}

\section{Conclusion}

\begin{frame}[label=future]{What Comes Next}
  \note[item]{
    Ce travail ouvre plusieurs perspectives de recherche.
  }

  \note[item]{
    D'abord, vérifier des optimisations plus complexes,
    telles que la propagation de constantes, ou la defunctionalisation selective.
  }

  \note[item]{
    Ensuite, étendre l'infrastructure pour composer des passes de compilation
    à travers différents langages et niveaux de représentations intermédiaires.
  }

  \note[item]{
    À terme, nous envisageons d'intégrer directement ce cadre dans CompCert
    et d'étendre la logique aux programmes concurrents avec gestionnaires d'effets.
  }

  \note[item]{
    Une autre question de voir si notre infrastructure peut etre
    integrer progressivement dans le pipeline de CompCert.

  }
  \note[item]{
    L'objectif à long terme est de travailler
    pour un compilateur optimisant entièrement vérifié
    pour un langage complexe tel que OCaml ou Rust.
  }

\end{frame}


\begin{frame}[label=thanks]{Thank you}
  \centering
  {\Large Questions?}
\end{frame}

\end{document}

% Local Variables:
% TeX-engine: luatex
% End:
